\section{正则理想}

设 A 是一个 k-代数。回忆（III, §1, No.~2, 第 430 页），A 的左理想是 A 的这样一个 k-子模 \(\mathfrak{a}\)：关系 \(a \in A\) 与 \(x \in \mathfrak{a}\) 蕴含 \(ax \in \mathfrak{a}\)。我们同样定义右理想与双边理想的概念。若 A 是一个 k-代数而 \(\mathfrak{a}\) 是 A 的一个双边理想，则在过渡到商时，A 中的乘法在 k-模 A/\(\mathfrak{a}\) 上定义一个 k-代数结构（出处同上）。

若 A 的一个左理想按包含关系是 A 的真左理想所成集合的极大元素，则称它是极大的。

定义 1. —— 设 A 是一个 k-代数，\(\mathfrak{a}\) 是 A 的一个左理想。A 的一个元素 u，若 au − a 属于 \(\mathfrak{a}\)（对每个 \(a \in A\)），则称为模 \(\mathfrak{a}\) 的右单位元。我们称理想 \(\mathfrak{a}\) 是正则的，如果存在模 \(\mathfrak{a}\) 的右单位元。

类似地，我们称 A 的右理想 \(\mathfrak{b}\) 是正则的，如果 A 有一个模 \(\mathfrak{b}\) 的左单位元，即一个元素 v，使得 \(va - a \in \mathfrak{b}\)（对每个 \(a \in A\)）。当一个理想是双边理想时，需要指明它作为左理想还是作为右理想是正则的。

当代数 A 有单位元时，A 的单位元素是模 \(\mathfrak{a}\) 的右单位元（对 A 的每个左理想 \(\mathfrak{a}\)）；这时，A 的每个左（或右）理想都是正则的。另一方面，若 A 是一个所有元素都是幂零元素的 k-代数，则 A 是 A 的唯一的正则（左或右）理想。

设 \(\mathfrak{a}\) 是 A 的一个正则左理想，\(u\) 是模 \(\mathfrak{a}\) 的一个右单位元。若 \(\mathfrak{b}\) 是 A 的一个包含 \(\mathfrak{a}\) 的左理想，则 \(u\) 是模 \(\mathfrak{a}\) 的一个右单位元，从而 \(\mathfrak{b}\) 是正则的。于是，A 的既极大又正则的左理想是 A 的正则真左理想所成集合的极大元素。此外，A 的一个左理想 \(\mathfrak{b}\)，若它包含 \(\mathfrak{a}\)，则是真的，当且仅当 \(u\) 不属于 \(\mathfrak{b}\)。因此，真左理想 \(\mathfrak{b}\)（A 的包含 \(\mathfrak{a}\) 的真左理想）所成之集按包含关系是一个归纳集。集合论 III, §2, No.~4, 第 154 页的定理 2 应用于这个集合就给出下面的结果。

命题 1. —— A 的每个正则真左（相应地，右）理想都包含于一个正则极大左（相应地，右）理想之中。

命题 2. —— 设 A 是一个交换 k-代数。A 的理想 \(\mathfrak{a}\) 是正则极大理想，当且仅当伪环 A/ \(\mathfrak{a}\)（I, §8, No.~1, 第 98 页）是一个体。

A 的元素 u 是模 \(\mathfrak{a}\) 的（右或左）单位元，当且仅当 u 在 A/\(\mathfrak{a}\) 中的典范像是代数 A/\(\mathfrak{a}\) 的单位元素。因此，理想 \(\mathfrak{a}\) 是正则的，当且仅当代数 A/\(\mathfrak{a}\) 有单位元。假设这些条件成立。

映射 \(\mathfrak{b} \mapsto \mathfrak{b}/\mathfrak{a}\) 是从代数 A 的包含 \(\mathfrak{a}\) 的理想所成之集到代数 A/\(\mathfrak{a}\) 的理想所成之集上的一个双射。由于这个代数有单位元，这些就是环 A/\(\mathfrak{a}\) 的理想。最后，由 I, §9, No.~1, 第 115 页的定理 1，环 A/\(\mathfrak{a}\) 是一个体，当且仅当它非零且其仅有的理想是 0 与它自身。命题得证。

例. —— 1) 设 V 是交换域 K 上的一个无限维向量空间。设 \(\operatorname{End}_{\mathrm{K}}^{\mathrm{f}}(\mathrm{V})\) 是 \(\operatorname{End}_{\mathrm{K}}(\mathrm{V})\) 的由有限秩自同态组成的 K-子代数。设 W 是 V 的一个向量子空间，并设 \(\mathfrak{a}_{W}\) 是 \(\operatorname{End}_{\mathrm{K}}^{\mathrm{f}}(\mathrm{V})\) 中核包含 W 的元素所成之集；它是 \(\operatorname{End}_{\mathrm{K}}^{\mathrm{f}}(\mathrm{V})\) 的一个左理想。\(\operatorname{End}_{\mathrm{K}}^{\mathrm{f}}(\mathrm{V})\) 的元素 u 是模 \(\mathfrak{a}_{W}\) 的右单位元，当且仅当我们有 \(u(x)=x\)（对每个 \(x\in W\)）。这样的元素 u 存在，也就是说 \(\mathfrak{a}_{W}\) 是正则的，当且仅当 W 是有限维的。理想 \(\mathfrak{a}_{W}\) 既极大又正则，当且仅当 W 的维数为 1。

*2) 设 \(\mathrm{T}\) 是一个局部紧空间，并设 \(\mathcal{C}_0(\mathrm{T})\) 是交换 \(\mathbf{C}\)-代数，它由从 \(\mathrm{T}\) 到 \(\mathbf{C}\) 的在无穷远处趋于 0 的连续映射组成（Gen.~Top., X, §4, No.~4, 第 316 页）；它有单位元，当且仅当 \(\mathrm{T}\) 是紧的。设 \(\mathrm{F}\) 是 \(\mathrm{T}\) 的一个闭子集，并设 \(\mathfrak{a}_{\mathrm{F}}\) 是 \(\mathcal{C}_0(\mathrm{T})\) 中在 F 上的限制为零函数的元素所成之集；它是 \(\mathcal{C}_{0}(T)\) 的一个理想。\(\mathcal{C}_{0}(T)\) 的元素 u 是模 \(\mathfrak{a}_{F}\) 的单位元，当且仅当我们有 \(u(t)=1\)（对每个 \(t\in F\)）。这样的元素 u 存在，也就是说 \(\mathfrak{a}_{F}\) 是正则的，当且仅当 F 是紧的。映射 \(t\mapsto \mathfrak{a}_{\{t\}}\) 是从 T 到 \(\mathcal{C}_{0}(T)\) 的正则极大理想所成之集上的一个双射（TS, I, §3, n° 2, 第 32 页，推论 1）。假设 T 不是紧的，并用 \(\mathfrak{a}\) 表示 \(\mathcal{C}_{0}(T)\) 中由具有紧支集的函数组成的子集；则 \(\mathfrak{a}\) 是 \(\mathcal{C}_{0}(T)\) 的一个理想，它不包含于 \(\mathcal{C}_{0}(T)\) 的任何正则理想之中。

\begin{enumerate}
\def\labelenumi{\arabic{enumi})}
\setcounter{enumi}{2}
\tightlist
\item
  设 \(L^{1}(\mathbf{R})\) 是局部紧群 R 的卷积代数。回忆（参见 Int., VIII, §4, No.~5, 第 38 页），\(L^{1}(\mathbf{R})\) 是 R 上关于勒贝格测度可积的函数的类所成的空间。两个函数 f 与 g 的类的乘积是由下述公式定义的函数 \(f * g\) 的类
\end{enumerate}

\[
(f * g) (s) = \int_ {- \infty} ^ {+ \infty} f (t) g (s - t) d t
\]

对几乎所有的 \(s \in \mathbf{R}\) 成立。代数 \(\mathrm{L}^1(\mathbf{R})\) 没有单位元。对任何 \(a\)（\(\mathbf{R}\) 中的元素），用 \(\mathfrak{m}_a\) 表示满足下述条件的元素 \(f\)（\(\mathrm{L}^1(\mathbf{R})\) 的元素）所成之集

\[
\int_ {- \infty} ^ {+ \infty} f (t) e ^ {- i a t} d t = 0.
\]

由 TS, II, §3, n° 2, 第 252 页，定理 1，映射 \(a \mapsto \mathfrak{m}_a\) 是从 \(\mathbf{R}\) 到代数 \(\mathrm{L}^1(\mathbf{R})\) 的正则极大理想所成之集上的一个双射。*

